% Appendix F — Ethics Statement

\subsection{Ethics Statement}
\label{app:ethics}

\paragraph{No harm to live systems.}
All proof-of-concept demonstrations were conducted on local devnet
infrastructure with no real assets at risk. The PEO exploit
(\S\ref{sec:peo}) uses a locally deployed, faithful reimplementation of the
mev-commit contract architecture on Foundry/Anvil. The PEFO structural
demonstration (\S\ref{sec:pefo-instance}) runs on an Astria local devnet
with synthetic rollup namespaces. No mainnet transactions were submitted and
no real funds were accessed or put at risk at any point.

\paragraph{Responsible disclosure.}
We followed responsible disclosure norms~\cite{iso29147}. Structural findings
were communicated to the affected shared sequencer and preconfirmation protocol
vendors prior to submission, with a 90-day window for response. All disclosure
attempts are timestamped via OpenTimestamps (Bitcoin calendar anchor) and
committed to a public artifact repository~\cite{artifact-repo}, providing
tamper-evident evidence of the disclosure timeline.

\paragraph{No vulnerability exploitation.}
The paper identifies structural properties of sequencing commitment protocols,
not exploitable vulnerabilities in specific deployed contracts. The PEFO
structural conditions (C1/C2/C3) are necessary consequences of the
lazy-sequencer architecture, documented in vendor publications. No
undisclosed zero-day vulnerability was identified or exploited.

\paragraph{Dual-use considerations.}
The PoC scripts in Appendix~\ref{app:poc} are designed for research
reproducibility. They require local deployment of test contracts and cannot
be directly applied to live systems without substantial adaptation. We have
not provided such adaptation and have no intention of doing so.

\paragraph{Survey and measurement.}
The $N=30$ application survey (Appendix~\ref{app:survey}) used only publicly
available information: deployed contract addresses, public documentation, and
on-chain data. No private or non-public information was accessed. The
12-month counterfactual replay uses publicly available block data from
production RPC endpoints; no proprietary MEV infrastructure was accessed.
